\section{System model}
\label{sec:model}

This section states the model compactly; the full parameter tables with units and sources are in
\texttt{docs/model/routing-optimisation.md} (decision 0014), and the implemented variant in
\texttt{docs/model/network-scheduler.md} (decisions 0024, 0025).

\paragraph{Sites and hours.} A network consists of sites $s \in \Sites$ (houses, hotels with spa,
indoor pools). Time runs in hourly steps $t$, local time Europe/Zurich. Site $s$ has $N_s$ servers,
recovers $a_s$ kWh of heat per server-hour ($a_s = 0.8$ in the case), has a minimum load share
$\mu_s$ while on, and a lossless tank of usable capacity $B_s$ kWh. Its gross heat demand (space
heating and hot water) in hour $t$ is $D_{s,t}$ kWh.

\paragraph{Realised heat balance.} Whatever a strategy decides, every hour is settled by one shared
function (\texttt{HeatBalance.Step}) on the \emph{actual} demand: server heat $a_s u_{s,t}$ first
covers demand, then charges the tank, and only the rest is dumped; the tank discharges for the
remaining demand, and the backup heater covers what is still missing. There is no second heat
model: the year simulation, the network simulation and the tests share this function.

\paragraph{Work.} We use two job types (decision 0020). \emph{Type~1} is paid batch LLM work
(for example summarising public documents); it arrives in \emph{cohorts} $c \in \Coh$, one per
origin site $o_c$ and day, at local 00:00, with an exclusive deadline $d_c$ \qty{24}{h} later and a
remaining amount $W_c$ server-hours. Its volume in our experiments is a labelled scenario: a share
of the origin's yearly heat-led server-hours, because no sourced arrival rate exists. Type~1 may
run only at eligible sites $E_c \subseteq \Sites$, which the data-locality rules (decision 0013)
determine before any optimisation; no price or heat term can override them. \emph{Type~2} is
\emph{filler}: research simulation with checkpoints and no deadline, available up to
$F_{s,t} = \varphi N_s$ server-hours per hour. $\varphi = 1$ is the case-permitted ``unlimited
jobs''; lower values model a shortage of sellable work. Work is divisible into small units, and
servers can run at part load, so loads are continuous.

\paragraph{Forecasts.} Plans for future hours use a demand forecast $\hat D_{s,t}$. By default it is
perfect foresight (the actual 2025 demand) and labelled as such. A synthetic forecast error follows
an AR(1) process in the lead time $k$:
\begin{equation}
  e_k = \rho\, e_{k-1} + \sigma_k \xi_k,\qquad \sigma_k = \sigma_{24}\sqrt{k/24},\qquad
  \xi_k \sim \mathcal{N}(0,1),\quad e_0 = 0,
  \label{eq:ar1}
\end{equation}
applied as a relative error to the demand and seeded per site and hour, so runs are reproducible.
The values of $\sigma_{24}$ and $\rho$ are \texttt{TODO(source)}; we therefore show a sweep
($\rho = 0.8$ is a placeholder).

\paragraph{Planning model.} At hour $\tau$ the scheduler plans the window
$\Win = \{\tau, \dots, \tau + H\}$ with $H = 48$ and commits only hour $\tau$ (rolling horizon,
decision 0014 \S6.2). It knows the realised tank content $S^0_s$ and the cohorts that have arrived
but are not finished; it does not know future arrivals. All variables are continuous and
non-negative: total load $u_{s,t}$, Type~1 work $x_{c,s,t}$ (defined only for $s \in E_c$ and
$t < d_c$), filler $f_{s,t}$, heat used directly $q^{\mathrm{dir}}_{s,t}$, tank charge
$q^{\mathrm{ch}}_{s,t}$ and discharge $q^{\mathrm{dis}}_{s,t}$, backup heat $b_{s,t}$, dumped heat
$z_{s,t}$, tank content $S_{s,t}$ and missed work $m_c$. For all $s \in \Sites$, $t \in \Win$:
\begin{align}
  u_{s,t} &= \textstyle\sum_{c} x_{c,s,t} + f_{s,t}, \qquad
  \mu_s N_s \le u_{s,t} \le N_s, \qquad f_{s,t} \le \varphi N_s, \label{eq:load}\\
  a_s u_{s,t} &= q^{\mathrm{dir}}_{s,t} + q^{\mathrm{ch}}_{s,t} + z_{s,t}, \label{eq:produce}\\
  q^{\mathrm{dir}}_{s,t} + q^{\mathrm{dis}}_{s,t} + b_{s,t} &= \hat D_{s,t}, \label{eq:demand}\\
  S_{s,t+1} &= S_{s,t} + q^{\mathrm{ch}}_{s,t} - q^{\mathrm{dis}}_{s,t}, \qquad
  0 \le S_{s,t} \le B_s, \qquad S_{s,\tau} = S^0_s, \label{eq:tank}\\
  \textstyle\sum_{s,t} x_{c,s,t} + m_c &= W_c \quad (d_c \le \tau + H), \qquad
  \textstyle\sum_{s,t} x_{c,s,t} \le W_c \quad (d_c > \tau + H). \label{eq:work}
\end{align}
For the current hour $\hat D_{s,\tau}$ is the measured demand. Eligibility, deadlines and
locality are enforced by which $x_{c,s,t}$ exist, not by penalties. Work processed at $s \ne o_c$
counts as \emph{moved}. A uniform transfer delay $\delta$ (whole hours, default 0; decision 0030)
models friction under the assumption ``data staged at arrival'': the cohort's data starts copying to
the eligible sites when it arrives, so $x_{c,s,t}$ with $s \ne o_c$ exists only for
$\alpha_c + \delta \le t < d_c$, where $\alpha_c$ is the arrival hour. This is a lower bound on
transfer friction; $\delta \ge \qty{24}{h}$ means no moves. There is no distance model and no
transfer cost (\texttt{TODO(source)}: a sourced data volume per server-hour); runs report moved
server-hours instead. The objective is lexicographic and is given in Section~\ref{sec:scheduling}.

\paragraph{Why an LP.} The problem is a flow in a time-expanded network with nodes $(s,t)$: work
waits at a site until its deadline, moves to an eligible site, or is computed; heat is a second
commodity that is used, stored or dumped, coupled to work by \eqref{eq:produce}
(\texttt{docs/model/routing-optimisation.md}~\S6.1). With divisible work and part load it stays a
linear programme. Integer servers, indivisible or non-preemptive jobs, start-up costs or
all-or-nothing payments make it a MILP (future work, Section~\ref{sec:limits}).
